
\documentclass[addpoints]{exam}
\usepackage{amsmath}
\usepackage{amssymb}
\usepackage{color}
\usepackage{siunitx}

% \printanswers

\newcommand{\event}{Forces \& Newton's Laws}
\newcommand{\tournament}{}
\def\type{0}

\setlength\fillinlinelength{\textwidth}
\CorrectChoiceEmphasis{\color{red}\bfseries}
\SolutionEmphasis{\color{red}\bfseries}

\setlength\answerclearance{2pt}
\pagestyle{headandfoot}
\runningheadrule
\runningheader{\event}{\tournament}{\ifnum\type=1 Class Set Test \else Full Name:\hspace{1cm} \fi}
\runningfootrule
\runningfooter{}{Page \thepage\ of \numpages}{}

\begin{document}
\begin{center}
    {\huge Grade 11 Physics \\ \event
    \ifprintanswers \space Key
    \else \space \ifnum\type<2 Test \fi \ifnum\type=1 \textbf{(CLASS SET)} \fi
          \ifnum\type=2 Answer Sheet \fi
    \fi \\ \vspace{3mm}\tournament\vspace{1mm}}
\end{center}

\vspace{.25\textheight}

\begin{center}
    First Name: \ifprintanswers
    \underline{\hspace{3.5cm}}\textbf{KEY}\underline{\hspace{5.5cm}}\\\vspace{5mm}
    \else \underline{\hspace{10cm}}\\ \vspace{5mm} \fi
    Last Name: \underline{\hspace{10cm}}\\ \vspace{5mm}
\end{center}

\noindent\textbf{\underline{Directions:}}
\begin{itemize}
    \ifnum\type=1 \item \textbf{This test is a class set.} Do not write on this paper. \fi
    \item Show all work with units and free body diagrams where required.
    Use $g = \SI{9.8}{m/s^2}$.
    \item The test is designed to be completed in 75 minutes.
\end{itemize}
\begin{center}
\textit{For grading use only}\\
\multirowgradetable{1}[pages]
\end{center}

\newpage
\begin{questions}

\section*{Part A --- Multiple Choice (15 marks)}
\vspace{5mm}

\question[1] Newton's First Law states that an object at rest will remain at rest unless:
\begin{choices}
    \choice It has zero mass
    \correctchoice A net external force acts on it
    \choice It is moving at constant velocity
    \choice Gravity is absent
\end{choices}

\question[1] A \SI{5.0}{kg} object accelerates at \SI{3.0}{m/s^2}. The net force on it is:
\begin{choices}
    \choice \SI{1.7}{N}
    \correctchoice \SI{15}{N}
    \choice \SI{8.0}{N}
    \choice \SI{45}{N}
\end{choices}

\question[1] Newton's Third Law states that:
\begin{choices}
    \choice Action and reaction forces act on the same object
    \correctchoice For every action force, there is an equal and opposite reaction force on a different object
    \choice Forces always cancel in pairs
    \choice Heavier objects exert greater reaction forces
\end{choices}

\question[1] The weight of a \SI{60}{kg} person on Earth is:
\begin{choices}
    \choice \SI{60}{N}
    \choice \SI{6.1}{N}
    \correctchoice \SI{588}{N}
    \choice \SI{600}{kg}
\end{choices}

\question[1] An object is in equilibrium. Which statement must be true?
\begin{choices}
    \choice It is at rest
    \choice It is moving at constant speed
    \correctchoice The net force on it is zero
    \choice No forces act on it
\end{choices}

\question[1] A \SI{10}{kg} box sits on a surface with $\mu_k = 0.30$. The kinetic friction
force is:
\begin{choices}
    \choice \SI{3.0}{N}
    \correctchoice \SI{29.4}{N}
    \choice \SI{98}{N}
    \choice \SI{0.30}{N}
\end{choices}

\question[1] Static friction is \textbf{greater than} kinetic friction because:
\begin{choices}
    \choice Objects at rest have more inertia
    \correctchoice More force is needed to start motion than to maintain it
    \choice $\mu_s < \mu_k$ always
    \choice Static friction depends on velocity
\end{choices}

\question[1] A block is on a frictionless inclined plane at angle $\theta$. The component
of gravity along the incline (down the slope) is:
\begin{choices}
    \choice $mg\cos\theta$
    \correctchoice $mg\sin\theta$
    \choice $mg\tan\theta$
    \choice $mg$
\end{choices}

\question[1] A spring with $k = \SI{200}{N/m}$ is stretched \SI{0.15}{m}. The restoring
force is:
\begin{choices}
    \choice \SI{0.75}{N}
    \choice \SI{200}{N}
    \correctchoice \SI{30}{N}
    \choice \SI{1333}{N}
\end{choices}

\question[1] An elevator accelerates upward. The normal force on a passenger inside is:
\begin{choices}
    \choice Equal to their weight
    \correctchoice Greater than their weight
    \choice Less than their weight
    \choice Zero
\end{choices}

\question[1] Two blocks of mass \SI{3.0}{kg} and \SI{5.0}{kg} are connected by a rope
and pulled by a force of \SI{24}{N} (frictionless surface). The acceleration of the system is:
\begin{choices}
    \choice \SI{3.0}{m/s^2}
    \correctchoice \SI{3.0}{m/s^2}
    \choice \SI{8.0}{m/s^2}
    \choice \SI{4.8}{m/s^2}
\end{choices}

\question[1] The normal force on a box resting on a horizontal surface equals its weight
when:
\begin{choices}
    \choice The box is moving
    \correctchoice The surface is horizontal and no vertical forces other than gravity act
    \choice The coefficient of friction is zero
    \choice The box is accelerating horizontally
\end{choices}

\question[1] A \SI{2.0}{kg} object hangs from two ropes, each making \ang{30} with the
ceiling. The tension in each rope is:
\begin{choices}
    \choice \SI{9.8}{N}
    \correctchoice \SI{11.3}{N}
    \choice \SI{19.6}{N}
    \choice \SI{5.0}{N}
\end{choices}

\question[1] A net force of \SI{0}{N} means:
\begin{choices}
    \choice No forces act on the object
    \correctchoice All forces on the object are balanced
    \choice The object must be at rest
    \choice The object is decelerating
\end{choices}

\question[1] Hooke's Law states:
\begin{choices}
    \choice $F = ma$
    \correctchoice $F = -kx$
    \choice $F = \mu mg$
    \choice $F = mg\sin\theta$
\end{choices}


\section*{Part B --- Short Answer (33 marks)}

\question A \SI{12}{kg} box is pushed across a horizontal floor by a horizontal force of
\SI{60}{N}. The coefficient of kinetic friction is $\mu_k = 0.25$.
\begin{parts}
    \part[2] Draw a free body diagram showing all forces on the box.
    \part[2] Calculate the normal force and the friction force.
    \part[3] Calculate the acceleration of the box.
    \part[2] If the pushing force is removed, how far does the box slide before stopping,
    given it was moving at \SI{4.0}{m/s} when the force was removed?
\end{parts}
\begin{solution}[10cm]
\textbf{(a)} FBD: weight $mg$ down, normal force $F_N$ up, applied force \SI{60}{N} right,
friction $f_k$ left.

\textbf{(b)} Vertical equilibrium: $F_N = mg = 12 \times 9.8 = \SI{117.6}{N}$

$f_k = \mu_k F_N = 0.25 \times 117.6 = \mathbf{\SI{29.4}{N}}$

\textbf{(c)} $F_\text{net} = F_\text{app} - f_k = 60 - 29.4 = \SI{30.6}{N}$
\[
a = \frac{F_\text{net}}{m} = \frac{30.6}{12} = \mathbf{\SI{2.55}{m/s^2}}$
\]

\textbf{(d)} After force removed: $F_\text{net} = -f_k = -29.4\,\text{N}$;\quad
$a = \dfrac{-29.4}{12} = \SI{-2.45}{m/s^2}$

$v^2 = v_0^2 + 2a\Delta d \Rightarrow 0 = 16 + 2(-2.45)\Delta d$
\[
\Delta d = \frac{16}{4.90} = \mathbf{\SI{3.27}{m}}
\]
\end{solution}

\question A \SI{5.0}{kg} block rests on a frictionless inclined plane at \ang{35} to
the horizontal.
\begin{parts}
    \part[2] Draw a free body diagram with all forces resolved along and perpendicular
    to the incline.
    \part[2] Calculate the component of gravity along the incline (parallel to slope).
    \part[2] Calculate the normal force.
    \part[3] If the incline has $\mu_k = 0.20$ (block is sliding down), find the net
    force and acceleration of the block.
\end{parts}
\begin{solution}[10cm]
\textbf{(a)} FBD: weight $mg$ down (components: $mg\sin\theta$ down the slope,
$mg\cos\theta$ into the slope); normal $F_N$ perpendicular out of slope;
friction $f_k$ up the slope (opposing motion).

\textbf{(b)} $F_\parallel = mg\sin\theta = 5.0 \times 9.8 \times \sin\ang{35}
= 49 \times 0.574 = \mathbf{\SI{28.1}{N}}$ down the slope

\textbf{(c)} $F_N = mg\cos\theta = 49 \times \cos\ang{35} = 49 \times 0.819
= \mathbf{\SI{40.1}{N}}$

\textbf{(d)} $f_k = \mu_k F_N = 0.20 \times 40.1 = \SI{8.02}{N}$ (up slope)
\[
F_\text{net} = F_\parallel - f_k = 28.1 - 8.02 = \mathbf{\SI{20.1}{N} \text{ down slope}}
\]
\[
a = \frac{F_\text{net}}{m} = \frac{20.1}{5.0} = \mathbf{\SI{4.02}{m/s^2} \text{ down slope}}
\]
\end{solution}

\question In an Atwood machine, a \SI{4.0}{kg} mass ($m_1$) and a \SI{6.0}{kg} mass ($m_2$)
hang on either side of a frictionless, massless pulley.
\begin{parts}
    \part[3] Draw the free body diagrams for both masses. Write Newton's second law for
    each.
    \part[3] Find the acceleration of the system.
    \part[2] Find the tension in the rope.
\end{parts}
\begin{solution}[10cm]
\textbf{(a)} For $m_1$ (lighter, accelerates up): $T - m_1 g = m_1 a$

For $m_2$ (heavier, accelerates down): $m_2 g - T = m_2 a$

\textbf{(b)} Add the two equations:
\[
m_2 g - m_1 g = (m_1 + m_2)a \implies a = \frac{(m_2 - m_1)g}{m_1+m_2}
= \frac{(6.0-4.0)(9.8)}{10.0} = \frac{19.6}{10.0} = \mathbf{\SI{1.96}{m/s^2}}
\]

\textbf{(c)} From $m_1$'s equation:
\[
T = m_1(g + a) = 4.0(9.8 + 1.96) = 4.0 \times 11.76 = \mathbf{\SI{47.0}{N}}
\]
\end{solution}

\question A spring scale in an elevator reads \SI{650}{N} for a person whose true weight is
\SI{588}{N} (mass = \SI{60}{kg}).
\begin{parts}
    \part[2] Draw the FBD of the person and write Newton's second law.
    \part[2] Find the acceleration of the elevator. State its direction.
    \part[2] What would the scale read if the elevator were in free fall?
    \part[2] A spring with $k = \SI{500}{N/m}$ hangs from the elevator ceiling with a
    \SI{2.0}{kg} mass. How much does the spring stretch when the elevator accelerates
    upward at \SI{1.07}{m/s^2}$?
\end{parts}
\begin{solution}[9cm]
\textbf{(a)} FBD: weight \SI{588}{N} down, normal (scale reading) \SI{650}{N} up.
$F_N - mg = ma$

\textbf{(b)} $a = \dfrac{F_N - mg}{m} = \dfrac{650 - 588}{60} = \dfrac{62}{60}
= \mathbf{\SI{1.03}{m/s^2} \text{ upward}}$

\textbf{(c)} In free fall, $a = g$ downward $\Rightarrow F_N = m(g - g) = \mathbf{\SI{0}{N}}$
(apparent weightlessness)

\textbf{(d)} Net upward force on mass: $F_\text{spring} - mg = ma$
\[
F_\text{spring} = m(g+a) = 2.0(9.8 + 1.07) = 2.0 \times 10.87 = \SI{21.74}{N}
\]
\[
x = \frac{F}{k} = \frac{21.74}{500} = \mathbf{\SI{0.0435}{m} \approx \SI{4.35}{cm}}
\]
\end{solution}

\end{questions}
\end{document}
